\section{The radical and extension of scalars}\label{sec:consequences}

The construction has already supplied the two Ext vanishings and a
finite nonprojective Gorenstein-projective module. We verify the
remaining assertions of Theorem~\ref{thm:main}.

\begin{proposition}\label{prop:radical}
The algebra $\Lambda$ is noncommutative, satisfies
$\Lambda/\rad\Lambda\simeq\kk^8$, and has
$\rad(\Lambda)^4\ne0$.
\end{proposition}\nopagebreak
\begin{proof}
Put $J_T=\rad T$. Since $T=C\ltimes DC$ and $C/\rad C\simeq\kk^2$,
\[
 J_T=\rad C\oplus DC,\qquad T/J_T\simeq\kk^2.
\]
The ideal $J_T\otimes T+T\otimes J_T$ in $A=T\otimes T$ is nilpotent:
the two summand ideals commute, and each is nilpotent. Its quotient is
$\kk^4$, so it equals $J_A=\rad A$. In the triangular algebra, the ideal
\[
 J_\Lambda=\begin{pmatrix}J_A&0\\F&J_A\end{pmatrix}
\]
is nilpotent as well. Indeed, a product has at most one off-diagonal
factor, with powers of $J_A$ on its two sides; sufficiently long
products vanish. Its quotient is $\kk^4\oplus\kk^4$. A nilpotent ideal
with semisimple quotient is the radical, proving the split assertion
and the count of simple modules.

Let $\chi:T\to\kk$ be the character of $s$. There are surjective
algebra maps
\[
 \Lambda\longrightarrow A
 \xrightarrow{\id_T\otimes\chi}T\longrightarrow C.
\]
For a surjection of finite-dimensional algebras, the image of the
radical is the radical of the quotient: the image is nilpotent, and
the quotient by that image is a quotient of a semisimple algebra.
In $C$, the multiplication table gives
\[
 (ut)^2=(y+qx)^2=q(1+q)z\ne0.
\]
Since $u,t\in\rad C$, the fourth radical power of $C$ is nonzero,
and consequently so is that of $\Lambda$. Finally $eu=u$ and $ue=0$
in $C$, so $C$ and its covering algebra $\Lambda$ are noncommutative.
\end{proof}

\begin{proposition}\label{prop:basechange}
For every field extension $K/\kk$, the algebra
$\Lambda_K=K\otimes_\kk\Lambda$ and module $Z_K=K\otimes_\kk Z$
satisfy all conclusions of Theorem~\ref{thm:main} over $K$.
\end{proposition}\nopagebreak
\begin{proof}
For finite $\Lambda$-modules $M,N$, take a resolution of $M$ with
finite projective terms. In each degree, finite projectivity gives
\[
 K\otimes_\kk\Hom_\Lambda(P_i,N)
 \simeq\Hom_{\Lambda_K}(K\otimes_\kk P_i,K\otimes_\kk N).
\]
Flatness of $K$ over $\kk$ commutes with kernels and images, hence
\begin{equation}\label{eq:basechange}
 K\otimes_\kk\Ext_\Lambda^i(M,N)
 \simeq\Ext_{\Lambda_K}^i(M_K,N_K)\qquad(i\geq0).
\end{equation}
This proves both Ext vanishings. The totally acyclic complex defining
$Z$ remains exact after tensoring with $K$; the same termwise Hom
isomorphism, now with target $\Lambda$, shows that its dual complex
over $\Lambda_K$ is exact. Thus $Z_K$ is Gorenstein-projective.

Choose a finite projective surjection $P\twoheadrightarrow Z$ with
kernel $N$. It does not split, because $Z$ is nonprojective. Its class
in $\Ext_\Lambda^1(Z,N)$ is nonzero and stays nonzero under
\eqref{eq:basechange}, since field extension is faithful. Thus the
extended sequence does not split and $Z_K$ is nonprojective.

The radical ideal in Proposition~\ref{prop:radical} stays nilpotent
and has quotient $K^8$ after scalar extension, so it is the radical of
$\Lambda_K$. Its nonzero fourth power stays nonzero. This also proves
the split and radical assertions over $K$.
\end{proof}
